% TOP_PROBLEM: {"id": 10000086, "title": "When is the uniqueness threshold below one?", "category_id": 19, "category": {"id": 19, "name": "probability", "display_name": "Probability", "description": "Problems involving probability theory, stochastic processes, and random structures.", "slug": "probability", "order_index": 19, "created_at": "2026-07-31T15:26:25.671Z"}, "status": "partially_solved", "problem_number": "AMR-099-0086", "set_id": 13, "statusQualification": "Restores the original paper’s site-percolation convention and mentions its corresponding bond-percolation formulation."}
% FIRST_POSTED: 2026-10-01
% ENTRY_KIND: statement_only
% UnsolvedMath ID 10000086; source catalogue code AMR-099-0086
% !TeX program = lualatex
% Definition and sources only; this file is not an LLM solution attempt.
\documentclass[11pt,a4paper]{article}
\usepackage[margin=25mm]{geometry}
\usepackage{fontspec}
\setmainfont{lmroman10-regular.otf}[BoldFont=lmroman10-bold.otf,
 ItalicFont=lmroman10-italic.otf,BoldItalicFont=lmroman10-bolditalic.otf]
\usepackage{amsmath,amssymb,mathtools}
\usepackage{unicode-math}
\setmathfont{latinmodern-math.otf}
\AtBeginDocument{\let\setminus\smallsetminus}
\usepackage{xurl,hyperref}
\hypersetup{colorlinks=true,urlcolor=blue}
\setcounter{secnumdepth}{0}
\setlength{\parindent}{0pt}
\setlength{\parskip}{5pt}
\setlength{\emergencystretch}{3em}
\begin{document}
\section*{When is the uniqueness threshold below one?}\label{10000086}
{\small UnsolvedMath ID 10000086 · AMR-099-0086 · Probability · AMR Open Problem Lists}
\subsection{Definitions and mathematical statement}
Let $G$ be an infinite connected locally finite graph. In Bernoulli site percolation with parameter $p\in[0,1]$, independently retain each vertex with probability $p$ and retain every original edge whose two endpoints survive. Let $N_\infty(p)$ count the infinite connected components of the resulting induced subgraph. The uniqueness threshold is
\[
p_u(G)=\inf\{p\in[0,1]:\mathbb P_p(N_\infty(p)=1)=1\}.
\]
The set is nonempty because $G$ is connected and infinite and $p=1$ retains it completely.

Give general geometric or structural conditions on $G$ implying $p_u(G)<1$. The specific proposed criterion is: if $G$ is vertex-transitive and has one end, must $p_u(G)<1$?

Vertex transitivity means that every pair of vertices can be exchanged by some graph automorphism. Having one end means that for every finite vertex set $F$, the graph obtained by deleting $F$ has exactly one infinite connected component. Both hypotheses concern the deterministic graph before percolation.

A positive answer to the particular criterion must work for every locally finite one-ended transitive graph, without assuming planarity, finite presentation of a group, or a particular growth rate. The conclusion asks for some parameter strictly below one at which there is almost surely exactly one infinite cluster; it does not require every vertex to belong to that cluster. The source uses site percolation and also proposes the analogous nonplanar questions for bond percolation, where edges rather than vertices are retained independently.
\subsection{Short English statement}
Does having one end force uniqueness of the infinite percolation cluster at some parameter below one on every transitive graph?
\subsection{Sources}
\begin{itemize}
\item[S1] ULAM.AI and the UnsolvedMath Contributors, supplied problems.json, record 10000086 (AMR-099-0086), AMR Open Problem Lists. Catalogue identity and source transcription; CC BY 4.0.
\url{https://huggingface.co/datasets/ulamai/UnsolvedMath}.
\item[S2] Benjamini and Schramm - Percolation beyond Z\textasciicircum{}d with 1999 update, Question 3, PDF page 9, PDF page 9. Original source indexed by AMR-099-0086.
\url{https://emis.de/ft/41630#page=9}.
\item[S3] I. Benjamini and O. Schramm, Percolation beyond Z\textasciicircum{}d, many questions and a few answers, Questions 3 and 5; site-percolation convention in Section 2.
\url{https://emis.de/ft/41630}.
\end{itemize}
\end{document}
